%==============================================================
%  Figura 6.1 - Determinazione del PL raggiunto nella fase di
%  realizzazione della SRP/CS: estratto del processo iterativo
%  di progettazione
%  Adattamento in italiano della Figura 6.1 dell'IFA Report 2/2017
%  (stesso stile della Figura 5.5, di cui e' la prosecuzione)
%==============================================================
\begin{tikzpicture}[
  font=\normalsize,
  grigio/.style={fill=black!10, align=center, inner sep=4pt},
  numcella/.style={draw=black, fill=white, line width=0.8pt,
                   minimum width=0.7cm, inner sep=1pt},
  corpo/.style={draw=black, fill=black!10, line width=0.8pt, align=center,
                minimum width=10.4cm, text width=10.0cm, inner sep=3pt},
  fdline/.style={draw=black, line width=0.9pt},
  fdar/.style={fdline, -{Triangle[length=6pt,width=6pt]}}
]

%--------------------------------------------------------------
% Ingresso dalla determinazione del PLr e nodo di rientro
%--------------------------------------------------------------
\node[align=center] (ingresso) at (6.0,-0.35)
  {Dalla determinazione del {PL\textsubscript{r}}\\ {(Figura 5.5)}};
\draw[fdar] (6.0,-0.95) -- (6.0,-1.45);
\draw[fdline] (6.0,-1.6) circle (0.15);

%--------------------------------------------------------------
% Fase 4: realizzazione delle SF
%--------------------------------------------------------------
\node[numcella, minimum height=0.75cm, anchor=west] (num4) at (3.3,-2.4) {4};
\node[corpo, minimum height=0.75cm, anchor=west] (fase4) at (4.0,-2.4)
  {Realizzazione delle SF, identificazione degli SRP/CS};

%--------------------------------------------------------------
% Fase 5: valutazione del PL (due righe)
%--------------------------------------------------------------
\node[numcella, minimum height=2.05cm, anchor=west] (num5) at (3.3,-4.425) {5};
\node[corpo, minimum height=1.3cm, anchor=west] (fase5a) at (4.0,-4.05)
  {Valutazione del PL degli SRP/CS in termini di\\ categoria, \textit{MTTF}\textsubscript{D},
   \textit{DC}\textsubscript{avg}, CCF};
\node[corpo, minimum height=0.75cm, anchor=west] (fase5b) at (4.0,-5.075)
  {Software e guasti sistematici};

%--------------------------------------------------------------
% Percorso principale
%--------------------------------------------------------------
\draw[fdar] (6.0,-1.75) -- (6.0,-2.025);
\draw[fdar] (6.0,-2.775) -- (6.0,-3.4);
\draw[fdar] (6.0,-5.45) -- (6.0,-6.3);

%--------------------------------------------------------------
% "Per ciascuna SF" con graffa
%--------------------------------------------------------------
\node[grigio, minimum width=1.5cm, minimum height=1.6cm] (perogni) at (1.15,-3.74)
  {Per\\ ciascuna\\ SF};
\draw[fdline, decorate, decoration={brace, amplitude=6pt}] (2.60,-5.45) -- (2.60,-2.025);

%--------------------------------------------------------------
% Uscite
%--------------------------------------------------------------
\node[align=center, text width=5.0cm] at (6.0,-7.2)
  {Alla verifica e\\ validazione (V\&V)\\ (Figura 7.1)};

\draw[fdline] (15.25,-6.1) -- (15.25,-1.6);
\draw[fdline] (15.55,-6.1) -- (15.55,-1.6);
\draw[fdar] (15.55,-1.6) -- (6.15,-1.6);
\node[align=center, text width=3.6cm] at (15.4,-7.45)
  {Ritorno se la\\ V\&V non ha\\ esito positivo\\ (Figura 7.1)};

%--------------------------------------------------------------
% Cornice
%--------------------------------------------------------------
\draw[draw=black, line width=1pt]
  ([shift={(-0.45,0.45)}]current bounding box.north west) rectangle
  ([shift={(0.45,-0.45)}]current bounding box.south east);

\end{tikzpicture}
